\section{Monitoring framework and evaluation design}
\label{sec:framework}

Both studies share the context policies, the per-turn quality signals, the
streaming monitor and the evidence standards defined here
(Figure~\ref{fig:framework}); Table~\ref{tab:datasets} summarises the two
benchmarks before the details below. Full model identifiers, gateway aliases
and acquisition dates are in supplementary Table~S18.

\begin{table}[!htbp]
\centering
\caption{The two benchmarks at a glance. K/L: same-context repeats / reworded
instructions per turn; ``full'' metric coverage means CQ, CS, RS, LS, ES and
SS, subject to the structural limits of Section~\ref{sec:framework} (CQ/RS/ES
need a final answer).}
\label{tab:datasets}
\footnotesize
\setlength{\tabcolsep}{4pt}
\begin{tabular}{@{}>{\raggedright\arraybackslash}p{0.17\linewidth}
>{\raggedright\arraybackslash}p{0.40\linewidth}
>{\raggedright\arraybackslash}p{0.40\linewidth}@{}}
\toprule
 & Study 1 (Lost in Conversation) & Study 2 (CDR-Stream)\\
\midrule
Source & GSM8K math, sharded \citep{cobbe2021training,laban2025lost} &
purpose-built, solver-checked\\
Units & \nTasks{} tasks, 4--9 turns (median 6) & 30 dialogues/model, 100 turns
each\\
Domains & arithmetic word problems & budget, inventory, lab, milestones (4)\\
Difficulty & not controlled & 12-level ladder, calibrated per model\\
Scenarios & single ordered reveal & 5 (2 stationary null, 3 scheduled-change)\\
Models & GPT-4o-mini, Llama-4 Scout, Mistral-7B & GPT-4o-mini, Llama-3 70B,
Gemma 3 12B\\
$K$ / $L$ & \Krep{} / \Lpar{}, all models & 3 / 2, GPT-4o-mini only\\
Metric coverage & full, pre-final turns only (no gold yet) & full for
GPT-4o-mini; CQ only for the other two\\
\bottomrule
\end{tabular}
\end{table}

\begin{figure}[!htbp]
\centering
\begin{tikzpicture}[
  font=\small, >=Latex, node distance=4mm and 6mm,
  box/.style={draw, rounded corners=2pt, align=center, minimum height=9mm,
              inner sep=3pt, fill=#1},
  box/.default=white]
\node[box=gray!10, text width=27mm] (user) {User turn $u_t$\\(new information,\\question)};
\node[box=blue!8, text width=33mm, right=of user] (policy)
  {Context policy $\pi$\\[1pt]\footnotesize dialogue history $\cdot$ fresh prefix\\
   \footnotesize $\cdot$ verified state summary};
\node[box=gray!10, text width=24mm, right=of policy] (llm) {LLM\\primary reply\\+ $K{-}1$ samples\\+ $L$ rewordings};
\node[box=orange!10, text width=26mm, below=9mm of llm] (metrics)
  {Per-turn quality\\vector $\mathbf{q}_t$\\\footnotesize CQ CS RS LS ES SS};
\node[box=orange!10, text width=27mm, left=of metrics] (window)
  {Window means\\$\bar q_{m,t}$ over the last\\$W$ valid values};
\node[box=red!8, text width=30mm, left=of window] (detect)
  {Drift detectors\\\footnotesize CUSUM, PH, EWMA,\\\footnotesize DDM, ADWIN};
\node[box=green!10, text width=30mm, below=8mm of detect] (alarm)
  {Alarm time $\tau_m$\\\footnotesize compared with the change $c$};
\node[box=green!10, text width=27mm, right=of alarm] (null)
  {Threshold from\\stationary nulls\\\footnotesize false alarms $\le\alpha$};
\node[box=gray!5, text width=26mm, right=of null] (gold)
  {Exact gold $y_t$\\\footnotesize(evaluation only)};
\draw[->] (user) -- (policy);
\draw[->] (policy) -- (llm);
\draw[->] (llm) -- (metrics);
\draw[->] (metrics) -- (window);
\draw[->] (window) -- (detect);
\draw[->] (detect) -- (alarm);
\draw[->, dashed] (null) -- (alarm);
\draw[->, dashed] (gold) -- (null);
\draw[->, dashed] (gold.north) -- (metrics.south);
\draw[->, dotted] (llm.north) -- ++(0,4mm) -| node[pos=0.25, above, font=\footnotesize]
  {primary reply enters later history (dialogue history only)} (policy.north);
\end{tikzpicture}
\caption{The monitoring framework. A context policy decides what the model
sees at turn $t$. Responses are scored into a per-turn quality vector, smoothed
over a window and passed to drift detectors whose thresholds are calibrated on
stationary dialogues. Exact gold is needed for CQ, RS and ES and for
evaluating alarms; CS, LS and SS do not use it.}
\label{fig:framework}
\end{figure}

\subsection{Context policies}

At turn $t$ the user sends $u_t$ and the model returns a primary reply $r_t$
from a context $C_t$ set by a policy $\pi$:
\begin{align}
  \text{dialogue history:}\quad & C_t = (u_1, r_1, \dots, r_{t-1}, u_t),\\
  \text{fresh prefix:}\quad & C_t = (u_1, \dots, u_t),\\
  \text{verified state summary:}\quad & C_t = (s_{t-1}, u_t),
\end{align}
where $s_{t-1}$ is the exact task state after turn $t-1$, rendered by a
solver. Dialogue history versus fresh prefix estimates the \emph{net} effect of
keeping assistant replies; context length and message positions also differ,
so the effect is not attributed to the content of the replies alone. The
verified state summary shortens the context, removes earlier replies and
supplies a correct state at once, so it tests a mitigation, not a mechanism.
Every context is packed into the request with server-side memory, retrieval and
prompt rewriting disabled, which makes it reproducible from the stored records.

\paragraph{A worked example (CDR-Stream)}
Suppose turn 7 reports ``warehouse~B received 40 crates of part~X,'' and turn
8 asks ``how many crates of part~X does warehouse~B hold now?'' Under
\emph{dialogue history}, $C_8$ is the full transcript of turns 1--8, including
whatever the model itself answered at every earlier turn; if an earlier update
or an earlier reply was wrong, $C_8$ still contains it, unflagged, alongside
the new turn. Under the \emph{verified state summary}, $C_8$ replaces that
transcript with the solver's current ledger for warehouse~B (one line per
part, with its correct running count) plus the new turn only; the model never
sees turns 1--7 or its own earlier replies. The gold answer is computed by the
same solver from the same ledger, independently of what either arm is shown.
Section~\ref{sec:study2} gives the full generator and prompt.

\subsection{Per-turn quality signals}

Each primary reply comes with $K-1$ further samples from the same context and
$L$ replies to frozen rewordings of the instruction. RQEval \citep{rqeval2026}
computes six dimensions in $[0,1]$, used unchanged: accuracy (CQ, the primary
answer equals the gold), consistency (CS, pairwise agreement of the $K$
answers), robustness (RS, agreement of reworded replies with a correct primary
answer), logical coherence (LS, absence of NLI contradictions between
consecutive steps; \citealp{he2021deberta}), efficiency (ES, correctness
combined with token economy) and stability (SS, BERTScore similarity of the $K$
traces; \citealp{zhang2020bertscore}). CQ, RS and ES need a final answer and the
gold, CS needs at least two answers, and LS and SS need only reasoning text.
Only LS and SS can therefore serve as label-free warnings before a model is
ready to answer. Concretely, LS splits a reply's reasoning into consecutive
steps, scores each adjacent pair for contradiction with an NLI model, and
reports one minus the mean contradiction probability, so a reply that never
contradicts itself scores 1 regardless of whether its conclusion is right. SS
instead compares the $K$ sampled \emph{replies} to the same context pairwise by
BERTScore and reports their mean similarity, so near-duplicate phrasing across
samples raises SS even when the shared answer is wrong. Both are computed only
when the reasoning text itself is present and parseable; missing values are
never imputed and are reported with their support (supplementary
Section~S1.2 gives the exact splitting and scoring rules).

\subsection{The streaming monitor}

Let $q_{m,t}\in[0,1]$ be the value of metric $m$ at turn $t$, or missing.

\begin{definition}[Window mean]
Let $V_{m,t}$ be the set of the last $W$ turns $j\le t$ at which $q_{m,j}$ is
observed. Once $|V_{m,t}|=W$, the monitored statistic is
\begin{equation}
  \bar q_{m,t} = \frac{1}{W}\sum_{j\in V_{m,t}} q_{m,j},
  \qquad W\in\{10,20,30\},
  \label{eq:windowmean}
\end{equation}
with $W=20$ primary. Missing values neither enter nor advance the window.
\end{definition}

Five detectors, CUSUM, Page--Hinkley, EWMA, DDM and ADWIN, are assigned to
metrics by three strategies fixed in the funded project design: CUSUM on every
metric, a distribution-specific mapping, and ADWIN on every metric
(supplementary Table~S8).

\begin{definition}[Alarm, calibration and detection]
For a detector with threshold $h$ reading $\bar q_{m,1},\dots,\bar q_{m,t}$,
the alarm time in a stream of $T$ turns is
\begin{equation}
  \tau_m(h) = \min\{t\le T : \text{the detector alarms at } t\},
  \label{eq:alarm}
\end{equation}
with $\tau_m=\infty$ if no alarm occurs. The threshold $h^\ast$ is the most
sensitive grid value whose first-alarm rate on stationary streams,
$\widehat{P}_0(\tau_m(h)\le T)$, is at most $\alpha$; drift data are never used.
For a change at turn $c$, an alarm is a \emph{false alarm} if
$\tau_m(h^\ast)<c$ and a \emph{detection} if $c\le\tau_m(h^\ast)\le c+33$.
\end{definition}

Study~1 estimates $\widehat{P}_0$ from permutation nulls with family-wise
$\alpha=0.05$, Study~2 from its stationary dialogues. Detecting a scheduled
change within 33 turns is not yet a warning. We call a signal an \emph{early
warning} only if, in addition, it alarms before the accuracy (CQ) alarm in the
same dialogue, keeps the null calibration, and does both across models.

\subsection{A paired readiness measure}

Sharded tasks allow an answer-free paired signal. Let $a_t^{\mathrm{F}}$ and
$a_t^{\mathrm{D}}$ be 1 if the fresh-prefix and dialogue-history arms attempt an
answer at turn $t$, and 0 if they ask for more information. The
\emph{Cumulative Resolution Lag} of a task is
\begin{equation}
  \mathrm{CRL} = \frac{1}{|\mathcal{T}|}\sum_{t\in\mathcal{T}}
  \max\!\left(a_t^{\mathrm{F}}-a_t^{\mathrm{D}},\,0\right),
\end{equation}
over the pre-final turns $\mathcal{T}$ at which both arms responded validly. It
measures how persistently the history arm lags an information-matched
counterfactual in becoming ready to answer, and it was fixed as an auxiliary
measure before the Study~1 data were collected.

\subsection{Evidence standards}

Each result carries one of four labels: \emph{confirmatory} (a hypothesis fixed
before its outcomes were computed), \emph{prespecified secondary},
\emph{exploratory} (rules fixed before computation on data already seen) or
\emph{post hoc}. In Study~1 the whole protocol, including the analysis, was
frozen before data collection. In Study~2 the design, sample, difficulty
mappings, prompt and truncation policy were frozen before the main run; the
analysis protocol was frozen when a budget pause had left 25,039 of 54,000
calls complete, before any main-run metric was computed (dated timeline in
supplementary Table~S19). Raw responses are stored before
scoring and never overwritten, failed gates are reported rather than
relabelled, and every number, table and figure is regenerated by scripts from
hashed inputs.
